\documentclass[12pt]{article}
\usepackage[a4paper,top=2cm,bottom=2cm,left=3cm,right=3cm,marginparwidth=1.75cm]{geometry}

%%%%%%%%%%
%packages%
%%%%%%%%%%
%Spracheinstellung / language setting / replace "english" if needed
\usepackage[english]{babel}
\usepackage[absolute,overlay]{textpos}
\usepackage[giveninits=true,sorting=none,style=numeric,maxbibnames=5]{biblatex}
\usepackage{csquotes}
\addbibresource{Sources.bib}
\usepackage{amsmath}
\usepackage{graphicx}
\usepackage{hyperref}
\usepackage{setspace}
\usepackage{amssymb}
\usepackage{xcolor}

 
%--------------------------------------------%
%--------EDIT CONTENTS OF THIS BLOCK---------%
%------------Title Page Content--------------%
%--------------------------------------------%


%%% ENTER THE THESIS TITLE HERE%%%
\newcommand{\TITLE}{
\textbf{\begin{spacing}{1}
Reinforcement Learning for Optimal Setpoint Determination in Energy Utility Systems
\end{spacing}}}
%%% ENTER THE TYPE OF THESIS HERE %%%
%%% Bachelor Thesis or Master Thesis
\newcommand{\THESIS}{
Master Thesis
}
%%% ENTER YOUR STUDY PROGRAMME HERE%%%
\newcommand{\STUDYPROGRAMME}{
Master of Science in Artificial Intelligence
}
%%% ENTER YOUR NAME HERE %%%
\newcommand{\NAME}{
Ashish Thapa
}
%%% ENTER YOUR DATE OF BIRTH HERE %%%
\newcommand{\DOB}{
04 March 1999
}
%%% ENTER YOUR PLACE OF BIRTH HERE %%%
\newcommand{\POB}{
Kaski, Nepal
}
%%% HIER DIE EIGENE MATRIKELNUMMER EINTRAGEN %%%
\newcommand{\MATRICULATIONNUMBER}{
5008270
}
%%% ENTER FIRST REVIEWER HERE %%%
\newcommand{\FIRSTREVIEWER}{
Title and name of first reviewer
}
%%% ENTER SECOND REVIEWER HERE %%%
\newcommand{\SECONDREVIEWER}{
Title and name of second reviewer
}
%%% ENTER ADVISER HERE %%%
\newcommand{\ADVISER}{
Title and name of adviser (optional)
}
%%% ENTER (real) SUBMISSION DATE HERE%%%
\newcommand{\SUBMISSIONDATE}{
25 June 2026
}
%%% ENTER THE ABSTRACT IN ENGLISH HERE %%%
%%% Should not exceed half a page %%%
\newcommand{\ABSTRACT}{
Example summary in English.
}
%%% ENTER THE ABSTRACT IN GERMAN HERE %%%
%%% Should not exceed half a page %%%
\newcommand{\ZUSAMMENFASSUNG}{
Beispielzusammenfassung in Deutsch.
}
%%% ENTER THE (exact) TASK DESCRIPTION HERE %%%
\newcommand{\TASKDESCRIPTION}{
Energy utility systems are the backbone of modern energy infrastructure. Their complexity ranges from small, localized systems to vast, highly interconnected infrastructure networks. Most of them operate under dynamic conditions where multiple interacting variables influence operational efficiency, cost, reliability, and stability. Conventionally, optimal setpoints for such systems are determined using mathematical optimization. These approaches rely on accurate system models and are typically formulated as constrained optimization problems. While effective for well-defined and moderately sized systems, classical optimization becomes computationally expensive as system complexity increases. Furthermore, time-varying dynamics, uncertainty from the integration of renewable sources, and model-plant mismatches makes it challenging to use purely model-based approaches.

This thesis investigates whether reinforcement learning can find optimal operating setpoints for two case studies of different complexity: a battery energy storage system and a non-linear cookie production factory. For each system, the control task is formulated as a Markov decision process (MDP) with explicit state, action, and reward definitions aligned with operational objectives. The RL agent is trained to produce adaptive setpoint policies that minimize operating cost while satisfying process constraints and preserving stable operation. The learned policies are then benchmarked against established mathematical optimization methods to determine where RL provides practical benefits and where its limitations remain.

The evaluation of the proposed method will be conducted through systematic simulation studies on two energy system case studies of varying complexity. Performance will be assessed using quantitative metrics including total operational cost, constraint violation frequency, and computational time. The results will provide a structured comparison between reinforcement learning and classical optimization and will determine under which conditions RL offers practical advantages for real-time setpoint optimization in complex energy utility systems. 
}
\begin{document}


%--------------------------------------------%
%--------FOLGENDEN BLOCK NICHT ÄNDERN--------%
%--------------------------------------------%


\pagenumbering{gobble}
\begin{center}
	{\Huge \TITLE}
\end{center}
\begin{textblock*}{\paperwidth}(0cm,8cm) % {block width} (coords) 
	\begin{center}
		{\Huge \textbf{\THESIS} \\ \textbf{\large \STUDYPROGRAMME}} \\[2em]
		{\Large submitted by} \\[2em]
		{\Huge \NAME} \\[0.5em]
		{\large \MATRICULATIONNUMBER} \\[2em]
		{\Large born on\DOB in\POB} \\[4em]
		{\Large AI-Lab}\\[0.5em]
		{\Large Chair of Applied Mathematics} \\[0.5em]
		{\Large Brandenburg University of Technology Cottbus-Senftenberg}
	\end{center}
\end{textblock*}
\begin{textblock*}{\paperwidth}(2.45cm,21cm) % {block width} (coords) 
	\begin{table}[!ht]
		\begin{tabular}{ll}
			{\Large Reviewer:}        & {\Large \FIRSTREVIEWER}  \\[1em]
			{\Large Reviewer:}        & {\Large \SECONDREVIEWER} \\[1em]
			{\Large Adviser:}         & {\Large \ADVISER}        \\[2em]
			{\Large Submission date:} & {\Large \SUBMISSIONDATE}
		\end{tabular}
	\end{table}
\end{textblock*}


\newpage\null\thispagestyle{empty}\newpage

~
\begin{textblock*}{\textwidth}(2.45cm,5cm) % {block width} (coords) 
	\section*{\underline{Declaration of independence}}
	\vspace{1cm}
	I hereby declare that I have written this thesis independently, without unauthorised help and exclusively using the sources indicated. The ideas taken directly or indirectly from external sources (including electronic sources) are labelled as such without exception. The sources used verbatim and in terms of content have been cited in accordance with the recognised rules of scientific work. This work has not yet been submitted for any other review or published elsewhere.\\[0.5cm]
	\noindent
	I agree that the thesis will be checked for plagiarism with the help of a plagiarism detection service.\\[2cm]
	\noindent
	{Cottbus, \SUBMISSIONDATE}{\hspace*{2.725cm}\underline{~~~~~~~~~~~~~~~~~~~~~~~~~~~~~~~~~~~~~~~~~~~~~~~~~~}} \\
	{\hspace*{7.75cm} Signature of the author}
\end{textblock*}

\newpage\null\thispagestyle{empty}\newpage
\section*{\underline{Task description}}
\vspace{0.5cm}
\TASKDESCRIPTION

\newpage\null\thispagestyle{empty}\newpage
\section*{\underline{Abstract}}
\vspace{0.5cm}
\ABSTRACT
\section*{\underline{Zusammenfassung}}
\vspace{0.5cm}
\ZUSAMMENFASSUNG


\newpage\null\thispagestyle{empty}\newpage
\tableofcontents

\newpage\null\thispagestyle{empty}\newpage
\pagenumbering{arabic}



%---------------------------------------%
%--------START EDITING FROM HERE--------%
%---------------------------------------%



\section{Introduction}
\label{Section1}
\textcolor{blue}{Introduces the research problem, motivation, and objectives of the thesis.}

% This is a chapter (see Sec.\ \ref{Section1}) with a figure (see Fig.\ \ref{IMG:ANN}).

% \begin{figure}[!h]
% \centering
% \includegraphics[scale=0.5]{TESTIMAGE}
% \caption{\label{IMG:ANN} Feedforward neural network. Source: \cite{schneidereit2022computational}}
% \end{figure}

% \subsection{This is a subsection}

% This is a subsection with a table (see Tab.\ \ref{TAB:NEURON}).

% \begin{table}[!h]
% \centering
% \begin{tabular}{c|c}
% Layer & Number of neurons \\\hline
% 1 & 5 \\
% 2 & 13 \\
% 3 & 10
% \end{tabular}
% \caption{\label{TAB:NEURON} Number of neurons in each layer.}
% \end{table}

% \subsubsection{This is a subsubsection}

% This is a subsubsection with an equation (see Eq.\ \eqref{EQ:SUM})

% \begin{equation}
% N(x)=\sum_{i=1}^5 x^2
% \label{EQ:SUM}
% \end{equation}

% \paragraph{This is a paragraph.}

% This is a paragraph.

\newpage
% is it better to have a section for Mathematical Optimization and Reinforcement Learning or to have them as subsections in the Theory and Background section? 
\section{Theory and Background}
\textcolor{blue}{Provides the fundamental concepts and background required to understand the systems, optimization methods, and reinforcement learning used in this work.}


Optimization is the process of selecting the best decision or solution, based on a defined objective, from a set of feasible alternatives while respecting given constraints. In energy systems, this typically involves determining control actions that minimize operational costs or energy consumption under physical and operational constraints.

Mathematically, optimization is formulated as:
\begin{equation}
\min_{x \in \Omega} f(x) \quad \text{s.t.} \quad h(x) = 0,\; g(x) \le 0,
\end{equation}
where $x \in \Omega \subseteq \mathbb{R}^n$ is the decision variable, $f : \mathbb{R}^n \to \mathbb{R}$ the objective function, $h : \mathbb{R}^n \to \mathbb{R}^p$ the equality constraints, and $g : \mathbb{R}^n \to \mathbb{R}^q$ the inequality constraints. The goal is to find the optimal solution $x^\star$ that achieves the minimum or maximum value $f(x^\star)$.

In practice, optimization problems can be addressed using different methods depending on the availability of accurate system models, the presence of uncertainty, and the complexity of decision-making.

In this thesis, two approaches are considered in the context of energy systems:
\begin{itemize}
    \item Mathematical optimization, which relies on explicit system models and deterministic formulations, and is used as a baseline for performance comparison;
    \item Reinforcement learning, which learns decision-making policies through interaction with the environment and constitutes the main approach of this work.
\end{itemize}

The following sections introduce the fundamental principles of both approaches.


\subsection{Mathematical Optimization}
\textcolor{blue}{Presents traditional optimization methods used for determining optimal system setpoints.}

Mathematical optimization uses explicit models of system behavior, constraints, and objectives to solve decision problems. In energy systems, it is commonly used to compute optimal operating strategies by minimizing costs or maximizing efficiency while respecting physical and operational limits. In this thesis, mathematical optimization is used as a baseline method for performance comparison, rather than being further developed as the main approach.

An optimization problem is typically formulated using a set of decision variables, an objective function, and a set of constraints. The decision variables represent the control actions to be determined, such as the charging and discharging of a battery or production levels in a factory. The objective function defines the performance criterion to be optimized, commonly expressed as cost minimization or profit maximization. Constraints capture the physical and operational limits of the system, including capacity bounds, energy balance equations, and technical restrictions.

Depending on the problem structure, optimization problems can be classified as linear, nonlinear, or mixed-integer. In energy system applications, mixed-integer linear programming (MILP) is widely used, as it enables the modeling of both continuous and discrete decision variables while maintaining computational tractability. However, as system size and complexity increase, such formulations can become computationally demanding and require accurate system modeling.

A key distinction in optimization is between deterministic and stochastic formulations. In deterministic optimization, all parameters, such as electricity prices or demand, are assumed to be known in advance. In contrast, stochastic optimization explicitly accounts for uncertainty by modeling uncertain parameters as random variables. While stochastic approaches provide a more realistic representation, they often lead to increased computational complexity.

Energy system optimization problems are often time-coupled, meaning that decisions made at one time step influence future system states. For example, the state of charge of a battery depends on previous charging and discharging actions, which introduces inter-temporal dependencies into the optimization problem.

These problems are typically solved using dedicated optimization solvers, such as those designed for linear and mixed-integer programming which can efficiently compute optimal solutions for well-defined models.
However, their performance depends strongly on the accuracy of system models and assumptions, and solution times can increase significantly for large-scale or highly dynamic problems, limiting their applicability in real-time decision-making \cite{kyriakidis2024hybrid}.

These limitations motivate the use of reinforcement learning, which can learn control strategies directly from interaction with the environment without requiring explicit system models, aiming to minimize operational costs or maximize economic profit, and therefore constitutes the main approach of this thesis.

\subsection{Reinforcement Learning}
\textcolor{blue}{Introduces reinforcement learning as a data-driven approach for learning optimal control policies through interaction with the environment.}

Reinforcement learning (RL) is a data-driven approach to sequential decision-making, in which an agent learns to take optimal actions through interaction with an environment. Unlike mathematical optimization, RL does not require an explicit model of system dynamics. Instead, it learns control strategies based on observed state transitions and received rewards, making it well-suited for complex and uncertain environments.

In the context of energy systems, reinforcement learning can be used to determine control actions, such as battery charging and discharging, in order to minimize operational costs or maximize economic profit over time. By continuously interacting with the system, the agent improves its decision-making policy based on experience.

Formally, reinforcement learning problems are typically modeled as a Markov Decision Process (MDP), defined by the tuple 
\[
(\mathcal{S}, \mathcal{A}, P, R, \gamma),
\]
where:
\begin{itemize}
    \item $\mathcal{S}$ is the set of states representing the system,
    \item $\mathcal{A}$ is the set of actions,
    \item $P(s' \mid s,a)$ defines the transition dynamics,
    \item $R(s,a)$ is the reward function,
    \item $\gamma \in [0,1]$ is the discount factor.
\end{itemize}

At each time step, the agent observes the current state $s$, selects an action $a$, and transitions to a new state $s'$ while receiving a reward. The objective is to learn a policy $\pi(a \mid s)$ that maximizes the expected cumulative reward over time \cite{sutton2018reinforcement}.

In energy system applications, the underlying decision problem can often be formulated as a stochastic optimal control problem, where uncertainties arise from exogenous variables such as renewable generation and electricity prices. These uncertainties are commonly modeled as stochastic processes, and the resulting problem can be expressed as a finite-horizon MDP \cite{pilling2025rl}. The goal is typically to minimize expected operational costs over time while respecting system dynamics and constraints.

Classical approaches such as dynamic programming can be used to solve such MDPs through backward recursion. However, these methods suffer from the curse of dimensionality when the state space becomes large. This limitation makes exact solutions computationally intractable for realistic energy systems.

Reinforcement learning provides an alternative by learning approximate solutions directly from data. As highlighted in \cite{pilling2025rl}, RL methods are especially useful in high-dimensional settings where traditional optimization methods become infeasible.

Furthermore, advanced RL methods such as policy gradient and actor-critic algorithms are better suited for problems with continuous or high-dimensional action spaces, as they avoid exhaustive discretization of actions and instead learn parameterized policies \cite{pilling2025rl}. This is particularly important for energy systems, where control decisions are naturally continuous.

Despite these advantages, reinforcement learning introduces new challenges, including the need for sufficient training data, potential instability during learning, and sensitivity to reward design and state representation.

In this thesis, reinforcement learning is used as the primary approach to learn control policies for energy system operation under uncertainty. Its performance is evaluated against a mathematical optimization benchmark, highlighting the trade-offs between model-based and data-driven approaches.

\section{Use Cases}
\textcolor{blue}{Describes the two systems considered in this thesis and their operational challenges.}

\subsection{Battery Energy Storage System}
\textcolor{blue}{Describes the battery system used as the first case study.}

Battery energy storage systems (BESS) are electrochemical devices that store electrical energy for later use and provide grid services such as energy arbitrage, peak shaving, and short-term balancing. A typical BESS is characterized by its usable energy capacity, maximum charge/discharge power, round-trip efficiency, and state-of-charge (SOC) operating limits. Real-world operation must balance economic objectives (e.g., buying low / using stored energy when prices are high) with engineering constraints (SOC bounds, power limits, efficiency losses).

BESS presents several challenges. Decisions are inherently time-coupled because current charging/discharging changes future availability. In addition, uncertain exogenous signals (market prices) require policies that are robust to forecast errors, and hard feasibility constraints must be enforced at every time step. These properties make BESS an illustrative case study for comparing optimization-based and reinforcement-learning approaches to real-time setpoint determination in energy systems.

A battery energy storage system provides flexible, time-shiftable power by charging when electricity prices are low and discharging when prices are high to reduce grid procurement cost. In this thesis, the battery is modeled using the \texttt{BatteryEnv} environment: the agent's observation includes the state-of-charge (SOC), a normalized current price, a normalized time index, and a forecast window of upcoming prices; the action is a continuous charge/discharge command $a\in[-1,1]$ scaled to real power $P\in[-P_{\max},P_{\max}]$.

The SOC dynamics are modeled as
\[
\mathrm{SOC}_{t+1} = \mathrm{SOC}_t + \frac{\eta\,P_{\text{actual}}\,\Delta t}{E_{\max}},
\]
where $\eta$ denotes round-trip efficiency, $E_{\max}$ the usable energy capacity, and $\Delta t$ the time step. The system is subject to the following operational constraints:

\begin{align}
0 \le \mathrm{SOC}_t \le 1, \quad \forall t, \label{eq:soc_bounds} \\
-P_{\max} \le P_t \le P_{\max}, \quad \forall t, \label{eq:power_limits} \\
P_{\text{actual}} = \begin{cases}
\eta_{\text{ch}} \cdot P_t & \text{if } P_t > 0 \text{ (charging)} \\
\frac{P_t}{\eta_{\text{disch}}} & \text{if } P_t < 0 \text{ (discharging)}
\end{cases} \label{eq:efficiency}
\end{align}

where $P_{\max}$ is the maximum charge/discharge power, $\eta_{\text{ch}}$ and $\eta_{\text{disch}}$ are the charging and discharging efficiency coefficients respectively. Constraint~\eqref{eq:soc_bounds} ensures the battery operates within safe operating levels, Constraint~\eqref{eq:power_limits} enforces physical power limitations, and Constraint~\eqref{eq:efficiency} models energy losses during charge and discharge cycles. The environment enforces feasibility by dynamically limiting the commanded power at each step so that the SOC remains within bounds.

The goal is to reduce the cost of electricity from the grid, which is calculated as the price multiplied by the power used. The environment gives the negative of this cost as the reward. In this work, there is no on-site generation and no selling back to the grid. The battery charges from the grid when prices are low and discharges to serve local demand when prices are high. Main challenges include respecting SOC and power limits, accounting for efficiency losses, handling uncertain future prices, and managing the fact that actions now affect future energy availability. These features make the BESS a suitable test case for reinforcement learning, where a policy must balance using energy now against saving it for likely higher-price periods while always obeying feasibility limits.

\subsection{Cookie Production Factory}
\textcolor{blue}{Introduces the industrial production system considered as the second case study.}

\section{Methodology}
\textcolor{blue}{Describes how the optimization problem is formulated as a reinforcement learning problem.}

\subsection{Problem Formulation}
\textcolor{blue}{Formulates the system control problem as a reinforcement learning problem.}

\subsection{Reinforcement Learning Framework}
\textcolor{blue}{Explains the structure of the RL agent and learning algorithm.}

\subsection{Training Setup}
\textcolor{blue}{Describes the training environment, reward function, and hyperparameter configuration.}

\subsection{Evaluation Metrics}
\textcolor{blue}{Defines the metrics used to evaluate the performance of the RL agent and optimization baseline.}

\section{Results and Discussion}
\textcolor{blue}{Presents and analyzes the experimental results.}

\subsection{Battery Energy Storage System}

\subsection{Cookie Production Factory}

\subsection{Key Findings}

\section{Conclusion and Future Work}
\textcolor{blue}{Summarizes the main findings and outlines possible future work.}

%Bibliography: Sources are stored in the bibtex file "Sources.bib" and referenced in the text via the reference \cite{Abbreviation}.
%Important note: Google Scholar offers the "Cite" feature - the information there must always be checked! They are often not completely correct!

\printbibliography

\end{document}